\scanpage{113}

\margin{?}

lin. indép. dans $H^2(S_0,\mathbb Z/2)$.\ldots

\uline{Pb} Considérons une forme quadratique \inserted{lisse} à coeff. indéterminés, à coefficients dans $\mathbb Z[(\unsure{a_{ij}|c_{ij}})_{1\leq i\leq j\leq n}]_{2\delta^{\prime}((c_{ij}))}$

``Calculer'' sa classe de Stiefel-Whitney totale ?

\uline{NB} Le $H^1(S_0,\mathbb Z/2)$ \inserted{$\simeq H^0(S_0,G_m)/2$} est engendré par $\varepsilon=\partial(-1),\ \deleted{\unsure{\partial|\xi}}\underbrace{\partial(2)}_{\xi},\ \underbrace{\partial(\delta'(c_{ij}))}_{\unsure{\eta|y}=w_1(f)}$, dont les 2 premiers deviennent nuls sur une extension biquadratique $\deleted{\unsure{\mathbb Q|\mathrm Q|\mathbb Z}}\mathbb Z_2[\sqrt{-1},\sqrt{2}]$ -- a fortiori sur $\mathbb C$, le deuxième donne lieu à $w_1$. Il est \inserted{donc} clair que multiplicativement à partir de ces éléments, on ne peut exprimer \inserted{aucun des} $w_2,w_3$ \ldots{} etc.

% archive sheet number: 112
